% Expanded, input-ready Section III. No standalone preamble or PDF build.
\section{Controller Realization and Complete Derivations}
\label{app:controller}
\noindent\textbf{Correspondence: Section~\ref{sec:fpi_controller}.} This appendix preserves the extended model, intermediate steps and supporting interpretation behind the corresponding main-text section.


\label{detail:sec:fpi_controller}
We propose a decentralized, occupancy-based incremental Fuzzy PI controller
for bittide. Each node combines receiving-buffer error with its
update-to-update change to compute an frequency correction increment.
This section specifies the complete controller and its continuous-clock
interface, derives an implementation-consistent PI calibration, and
establishes the nonlinear properties used in
Section~\ref{detail:sec:btd_reference}. The physical plant, integer frame counts
and sampled controller state are kept distinct throughout the derivation.

\subsection{Continuous-clock realization of the proposed controller}

\label{detail:subsec:fpi_laws}
The plant is the current continuous-phase Callisto model within the bittide
framework~\cite{lall2024logicalSynchrony}. For constant
uncontrolled clock frequency $\omega_i^u$ and applied correction $c_i(t)$,
\begin{align}
 \dot\theta_i(t)&=\omega_i^u+c_i(t),\nonumber\\
 c_i(t)&=c_{i,k},\qquad t\in[t^a_{i,k},t^a_{i,k+1}),\nonumber\\
 \theta_i(t)&=\theta_i(t^a_{i,k})+
       (\omega_i^u+c_{i,k})(t-t^a_{i,k}).
 \label{detail:eq:fpi_continuous_plant}
\end{align}
The superscript $a$ denotes actuation, and $k$ indexes that node's updates.
Normal executions require $\omega_i^u+c_i(t)>0$, so the phase has a
monotone inverse. Physical phase is continuous at an update; its slope can change. Requested
correction $q_{i,k}$ is computed from the measured buffers and applied as
$c_{i,k}=\operatorname{sat}_{c_{\max}}(q_{i,k})$, with
$\operatorname{sat}_b(z)=\min\{b,\max\{-b,z\}\}$ and $c_{\max}>0$.
The proposed controller has jump updates inside a continuous physical plant;
it is not implemented as an ordinary differential equation for its accumulator.

Let $p>0$ be the polling period and $d_{\rm ctrl}\ge0$ the execution delay,
both measured in local-clock phase units. Measurement time $t^m_{i,k}$
and actuation time $t^a_{i,k}$ satisfy
\begin{align}
 \theta_i(t^m_{i,k})&=\theta_i(0)+k\,p,\nonumber\\
 \theta_i(t^a_{i,k})&=\theta_i(0)+k\,p+d_{\rm ctrl},\nonumber\\
 \int_{t^m_{i,k}}^{t^a_{i,k}}(\omega_i^u+c_i(s))\,ds&=d_{\rm ctrl}.
 \label{detail:eq:fpi_local_schedule}
\end{align}
The product $k\,p$ in these scheduling equations means the integer index
$k$ multiplied by the poll period $p$; it is not the proportional gain
$k_p$. With zero initial correction, $t^m_{i,0}=0$ and
$t^a_{i,0}=d_{\rm ctrl}/\omega_i^u$. Subsequent actuation intervals equal
$p/(\omega_i^u+c_{i,k})$. The delay is not automatically a constant
physical time, and its age must not be computed using the new command
that is applied only at the end of that age interval.

In normal operation, each integer source-phase crossing sends a frame.
With the progress guards open, frame conservation gives
\begin{equation}
 \beta_{j\to i}(t)=\beta_{j\to i}(0)+A^{\rm arr}_{ji}(t)-N_i(t),
 \label{detail:eq:fpi_count_balance}
\end{equation}
where $A^{\rm arr}_{ji}(t)$ counts arrivals during $(0,t]$, including
initial in-flight frames, and
$N_i(t)=\lfloor\theta_i(t)\rfloor-\lfloor\theta_i(0)\rfloor$ counts
local consumption. For a constant physical delay $\tau_{ji}\ge0$,
\begin{equation}
 A^{\rm arr}_{ji}(t)=\lfloor\theta_j(t-\tau_{ji})\rfloor
                  -\lfloor\theta_j(-\tau_{ji})\rfloor.
 \label{detail:eq:fpi_fixed_delay_arrivals}
\end{equation}
Thus arrival counts incorporate the declared negative-time history.

For constant physical link latency $\tau_{ji}$, receiving occupancy is
\begin{equation}
 \beta_{j\to i}(t)=\lambda^{\rm UGN}_{ji}
 +\lfloor\theta_j(t-\tau_{ji})\rfloor-\lfloor\theta_i(t)\rfloor.
 \label{detail:eq:fpd_continuous_fifo}
\end{equation}
Here $\lambda^{\rm UGN}_{ji}$ is the fixed integer offset in the absolute
phase-count convention. It is determined by occupancy and prehistory:
\begin{equation}
 \lambda^{\rm UGN}_{ji}=\beta_{j\to i}(0)
 -\lfloor\theta_j(-\tau_{ji})\rfloor+\lfloor\theta_i(0)\rfloor.
 \label{detail:eq:fpd_ugn_init}
\end{equation}
The relative-count logical marking is
$\lambda^{\rm LSN}_{ji}=\lambda^{\rm UGN}_{ji}
+\lfloor\theta_j(0)\rfloor-\lfloor\theta_i(0)\rfloor$;
the two coincide when the initial phase counts are equal.

For changing delivery times, let $s_{j,m}=\theta_j^{-1}(m)$ be the send
time of source tick $m$, and $a_{ji}(s)$ its causal ordered arrival map.
For the desired send-time delay profile $\tau_{ji}(s)$, the declared
running-maximum envelope is
\[
 a_{ji}(s)=\sup_{r\le s}\{r+\tau_{ji}(r)\},
\]
including the constant-delay negative-time history.
Define the delivered tick frontier
\begin{equation}
 \begin{aligned}
 M^{\rm arr}_{ji}(t)&=\max\{m\in\mathbb Z:a_{ji}(s_{j,m})\le t\},\\
 \beta_{j\to i}(t)&=\lambda^{\rm UGN}_{ji}
       +M^{\rm arr}_{ji}(t)-\lfloor\theta_i(t)\rfloor.
 \end{aligned}
 \label{detail:eq:fpd_ordered_arrivals}
\end{equation}
The negative-time history specifies initial in-flight frames. The declared
running-maximum arrival envelope preserves order and can release a batch
after a desired-delay decrease; it imposes no finite serialization rate.
At an upward-delay discontinuity, a supremum of continuous source times
need not itself have arrived. Counting delivered tick indices avoids
silently including such a boundary frame.

If a protection guard blocks logical progress, the runtime can retain a
separate continuous data phase $\psi_i$ with
$\dot\psi_i=g_i(t)\dot\theta_i$, where $g_i\in\{0,1\}$ is the progress
gate. Buffer counts then use $\psi_i$, whereas scheduling still uses
$\theta_i$. In the ungated normal regime, with matching initial phases, data phase and physical phase agree. Guard operation does not stop the oscillator.

Only for the reference proof do we replace integer occupancy by a
fluid counterpart. At constant delay, define
$\widetilde\beta_{ji}(t)=\widetilde\lambda_{ji}
+\theta_j(t-\tau_{ji})-\theta_i(t)$, where
$\widetilde\lambda_{ji}$ is an additive offset chosen for the specified
fluid initialization. Between command discontinuities,
\begin{equation}
 \dot{\widetilde\beta}_{ji}(t)
 =\omega_j^u+c_j(t-\tau_{ji})-\omega_i^u-c_i(t).
 \label{detail:eq:fpi_fluid_counterpart}
\end{equation}
This derivative belongs to the fluid counterpart, not the actual integer
occupancy. Replacing floor counts by phases does not automatically preserve
an arbitrary initialization; the compatible zero-latency marking used in
Section~\ref{detail:sec:btd_reference} is specified there.

\subsection{Notation and the measured error}

Node $i$ has incoming-neighbor set $\mathcal N_i^-$ and $d_i=|\mathcal N_i^-|>0$ incoming channels. All incoming sums below range over $j\in\mathcal N_i^-$. This section assumes fixed active
incoming sets: the runtime link-status mask is one on all configured inputs.
The Fuzzy percentage wrapper averages its configured inputs, so changing
link masks would require an explicit corresponding measurement policy.
At the $k$th local measurement,
$\beta_{ji,k}$ is the measured occupancy of channel $j\to i$, and
$\beta^*_{ji}$ is its fixed target. Use
\begin{equation}
 E_{i,k}=\sum_{j\in\mathcal N_i^-}(\beta_{ji,k}-\beta^*_{ji}).
 \label{detail:eq:fpi_error}
\end{equation}
The sign is actual minus target. If incoming buffers are too full, a
positive frequency correction increases local consumption. This explains
the positive PI gains with this error convention. The physical frequency is $\omega_i=\omega_i^u+c_i$, where $\omega_i^u$ is the constant uncontrolled frequency; $c_i$ is a correction, not the entire frequency.

\begin{center}
\begin{tabular}{@{}ll@{}}
\toprule
Symbol & Meaning \\
\midrule
$E_{i,k}$ & Sum error in frames \\
$e_{i,k}$ & Target-minus-measured error in percentage points \\
$I_{i,k}$ & PI baseline sum of sampled frame errors \\
$A_{i,k}$ & Fuzzy accumulator: total requested correction \\
$q_{i,k},c_{i,k}$ & Requested and applied frequency corrections \\
$c_{\max}$ & Magnitude limit on the applied correction \\
$\beta_{\max}$ & Common buffer capacity in frames \\
$\beta_{\max,ji}$ & Capacity of receiving channel $j\to i$ \\
$\chi_{ji}$ & Used outgoing credit in frames \\
$h$ & Fixed PI baseline integration scale \\
$X_{i,k},Y_{i,k}$ & Scaled Fuzzy error and sample-difference inputs \\
$F_0,\mathcal R_F$ & Linear counterpart and residual $F-F_0$ \\
\bottomrule
\end{tabular}
\end{center}

The event index $k$ is local to a node. The normalization factor $\gamma_i$ converts a frame-sum error into percentage points. The maximum correction $c_{\max}>0$ is the actuator limit; $q$ and $c$ remain distinct if clipping acts. Two nodes need not measure at the
same physical times. The subscript $ji$ abbreviates edge $j\to i$.
Suppressing $i$ below is a notational convenience, not a synchrony
assumption; in that notation $\gamma=\gamma_i$.

\subsubsection{Deriving the percentage-to-frame-error relation}
For common incoming capacity $\beta_{\max}$, define
\[
 \bar p_{i,k}=\frac1{d_i}\sum_j\frac{100\beta_{ji,k}}{\beta_{\max}},\qquad
 r_i=\frac1{d_i}\sum_j\frac{100\beta^*_{ji}}{\beta_{\max}}.
\]
Then subtraction and factoring give
\begin{align}
 e_{i,k}&=r_i-\bar p_{i,k}\\
 &=-\frac{100}{d_i\beta_{\max}}\sum_j(\beta_{ji,k}-\beta^*_{ji})
 =-\gamma_i E_{i,k},\qquad \gamma_i=\frac{100}{d_i\beta_{\max}}.
 \label{detail:eq:fpi_percent}
\end{align}
For midpoint targets this is the original absolute-percentage input with
$r_i=50$. For other targets the runtime can supply a centered percentage
$r_i+d_i^{-1}\sum_j100(\beta_{ji,k}-\beta^*_{ji})/\beta_{\max}$; subtracting it from
$r_i$ yields the same $e_{i,k}$. A mismatched setpoint instead adds a
constant bias and invalidates the scalar relation used for calibration.

With unequal capacities and a consistently matched target reference, the
percentage average gives
\[
 e_{i,k}=-\frac{100}{d_i}\sum_j
 \frac{\beta_{ji,k}-\beta^*_{ji}}{\beta_{\max,ji}}.
\]
Here $\beta_{\max,ji}$ is the capacity of receiving channel $j\to i$.
This is a weighted error and generally cannot be written as one normalization constant
times the unweighted PI error sum. Matching such a controller requires a
consistent measurement weighting or a different baseline.

\subsection{PI baseline: update order, units and the first call}

The occupancy-feedback PI baseline is grounded in the bittide analysis
of~\cite{lall2022resistance}. The implemented PI update is an
integrate-then-output recursion:
\begin{equation}
 I_k=I_{k-1}+E_k,\qquad
 q_k^{\rm PI}=k_pE_k+\kappa_I I_k,\qquad \kappa_I=k_i h.
 \label{detail:eq:fpi_native}
\end{equation}
The sum $I_k$ contains no physical time multiplier; $h$ is incorporated in
$\kappa_I$. In the present runtime,
\begin{equation}
 h=\frac{p}{\omega_{\rm base}}.
 \label{detail:eq:fpd_h}
\end{equation}
Here $p$ is the local polling period and $\omega_{\rm base}$ is the fixed implementation reference frequency. The scale $h$ stays fixed during the configured run. Local polling intervals measured
in physical time can vary as the local clock frequency changes; replacing $h$ by
those intervals would implement a different PI law.

Subtract the previous requested output:
\begin{align}
 q_k^{\rm PI}-q_{k-1}^{\rm PI}
 &=k_p(E_k-E_{k-1})+\kappa_I(I_k-I_{k-1})\\
 &=k_p(E_k-E_{k-1})+\kappa_I E_k.
 \label{detail:eq:fpi_pi_increment}
\end{align}
This is the correct increment to reproduce. It contains the present
$E_k$ in the integral term because integration precedes output.

At zero initial sum, the first call returns
\begin{equation}
 q_0^{\rm PI}=(k_p+\kappa_I)E_0.
 \label{detail:eq:fpd_first}
\end{equation}
No physical prehistory measurement is needed to express this algebra:
one may use the fictitious previous error $E_{-1}=0$ and previous requested
output $q_{-1}=0$. That bookkeeping choice is separate from deriving a
plant identity such as $x_k-x_{k-1}=H w_{k-1}$, which requires a genuine
elapsed plant interval.

\subsection{The proposed Fuzzy PI controller}

The proposed local incremental control law combines the present error
and its sample difference as follows:
\begin{align}
 X_k&=K_e e_k,&Y_k&=K_{\Delta e}(e_k-e_{k-1}),\\
 A_k&=A_{k-1}+K_uF(X_k,Y_k),&q_k^{\rm F}&=A_k.
 \label{detail:eq:fpi_fuzzy_law}
\end{align}
The gain signs used here are $K_e,K_{\Delta e}>0$ and $K_u<0$. Thus high occupancy
produces negative $e$, and a negative surface output produces a positive
correction increment. The quantity $e_k-e_{k-1}$ is not divided by elapsed
time; calling it a time derivative would obscure its sampling dependence.

\subsubsection{Memberships and rule weights}
The labels NL, NS, ZE, PS and PL mean negative large, negative small, zero, positive small and positive large, respectively. Each membership $\varphi_a(z)\in[0,1]$ is an activation degree; a singleton is a numerical rule consequent. Let $[z]_+=\max\{z,0\}$ and $\operatorname{clip}_{[0,1]}z$
denote projection onto $[0,1]$. The exact memberships can be written
\begin{align}
 \varphi_{\rm NL}(z)&=\operatorname{clip}_{[0,1]}\frac{-z-5}{5},&
 \varphi_{\rm NS}(z)&=[1-|z+5|/5]_+,\nonumber\\
 \varphi_{\rm ZE}(z)&=[1-|z|/5]_+,&
 \varphi_{\rm PS}(z)&=[1-|z-5|/5]_+,\nonumber\\
 \varphi_{\rm PL}(z)&=\operatorname{clip}_{[0,1]}\frac{z-5}{5}.
 \label{detail:eq:fpi_full_memberships}
\end{align}
Table~\ref{detail:tab:fpi_full_rules} gives all 25 consequent rules.
\begin{table}[t]
\caption{Rule consequents of the proposed bittide Fuzzy PI controller.}
\label{detail:tab:fpi_full_rules}
\centering
\begin{tabular}{@{}c|ccccc@{}}
\toprule
$X\backslash Y$ & NL & NS & ZE & PS & PL\\
\midrule
NL & NL & NL & NL & NS & ZE\\
NS & NL & NL & NS & ZE & PS\\
ZE & NL & NS & ZE & PS & PL\\
PS & NS & ZE & PS & PL & PL\\
PL & ZE & PS & PL & PL & PL\\
\bottomrule
\end{tabular}
\end{table}
Order the labels as NL, NS, ZE, PS and PL, with label indices
$a,b\in\{1,\ldots,5\}$. With singleton values
$(s_1,\ldots,s_5)=(-5,-2.5,0,2.5,5)$ and
$\ell(a,b)=\min\{5,\max\{1,a+b-3\}\}$, compute
\begin{equation}
 w_{ab}=\min\{\varphi_a(X),\varphi_b(Y)\},\qquad
 F(X,Y)=\frac{\sum_{a,b}w_{ab}s_{\ell(a,b)}}{\sum_{a,b}w_{ab}}.
 \label{detail:eq:fpi_rules}
\end{equation}
The membership partition gives at least one positive input membership on
each axis, hence a positive denominator at finite inputs. Individual rules
are summed even when they share a consequent. Replacing this accumulation
by maximum aggregation changes the map and its derivation.

The weighted-average property gives $|F|\le5$ and therefore
$|A_k-A_{k-1}|\le5|K_u|$. This is a per-update bound, not a bound on $A_k$
over an arbitrarily long execution. The raw Fuzzy module itself has no
final correction clipping.
Initialize $A_{i,-1}=0$ and $e_{i,-1}=0$. At each measurement, form
$e_{i,k}$ and $(X_{i,k},Y_{i,k})$, evaluate the rule weights and normalized
surface, increment $A_{i,k}$, and store $e_{i,k}$ as the previous error.
The requested correction is then limited and applied at $t^a_{i,k}$, and
held until the next actuation. In the normal-run policy considered here,
external clipping retains the raw accumulator and anti-windup is disabled.
The controller does not overwrite $A_{i,k}$ with the applied command.
The membership functions, rule consequents and gains stay fixed during a
run; the correction changes with the measured input pair. The law requires
neither an observer nor an FFP mode.


\subsection{Exact PI Calibration, Including Initialization}
\label{detail:subsec:fpi_calibration}

Define the globally linear counterpart
\begin{equation}
 F_0(X,Y)=\frac{X+Y}{2}.
 \label{detail:eq:fpi_f0}
\end{equation}
Replacing $F$ by this map retains the same controller accumulator.
Using $e=-\gamma E$ gives
\begin{align}
 A^0_k-A^0_{k-1}
 &=\frac{K_u}{2}\{K_e e_k+K_{\Delta e}(e_k-e_{k-1})\}\\
 &=-\frac{K_u\gamma K_e}{2}E_k
   -\frac{K_u\gamma K_{\Delta e}}2(E_k-E_{k-1}).
 \label{detail:eq:fpi_substitution}
\end{align}
Matching the two independent coefficients in~\eqref{detail:eq:fpi_pi_increment}
yields
\begin{equation}
 k_p=-\frac{K_u\gamma K_{\Delta e}}2,\qquad
 \kappa_I=-\frac{K_u\gamma K_e}2.
 \label{detail:eq:fpi_matching}
\end{equation}
Solving for two Fuzzy gains with chosen $K_e>0$ gives
\begin{equation}
 K_u=-\frac{2k_ih}{\gamma K_e},\qquad
 K_{\Delta e}=\frac{K_ek_p}{k_ih}.
 \label{detail:eq:fpi_design_gains}
\end{equation}
The formulas assume $k_i h>0$. Pure proportional control is a separate
limiting case and should not be inserted into this inverse mapping with
a zero denominator.

\begin{btdproposition}[Event-level linear-counterpart equivalence]
\label{detail:prop:fpi_equivalence}
For any measurement sequence, suppose~\eqref{detail:eq:fpi_matching} holds and
initialize
\begin{equation}
 A^0_{-1}=k_pE_{-1}+\kappa_I I_{-1},\qquad e_{-1}=-\gamma E_{-1}.
 \label{detail:eq:fpi_initialization}
\end{equation}
Then
$A^0_k=k_pE_k+\kappa_I I_k$ for every $k\ge0$.
\end{btdproposition}
\begin{proof}
Assume the relation holds before update $k$. Substitution gives
\begin{align*}
 A^0_k
 &=k_pE_{k-1}+\kappa_I I_{k-1}
   +k_p(E_k-E_{k-1})+\kappa_I E_k\\
 &=k_pE_k+\kappa_I(I_{k-1}+E_k)
 =k_pE_k+\kappa_I I_k.
\end{align*}
The stated initialization supplies the induction base, including $k=0$.
\end{proof}

\subsubsection{Why zero constructors work for a nonzero first error}
Take $I_{-1}=A^0_{-1}=e_{-1}=0$, corresponding to $E_{-1}=0$.
The first Fuzzy-reference input is
\[
 (X_0,Y_0)=(-\gamma K_eE_0,-\gamma K_{\Delta e}E_0),
\]
so
$A^0_0=(k_p+\kappa_I)E_0$, exactly as in~\eqref{detail:eq:fpd_first}.
If only Fuzzy is warmed up to $e_{-1}=e_0$ while $A^0_{-1}=0$, its
difference input vanishes and the output is merely $\kappa_I E_0$.
The missing $k_pE_0$ term can create a large apparent controller difference
unrelated to fuzzy inference.

\subsubsection{A numerical calibration from the present configuration}
For $d=2$, $\beta_{\max}=2048$, $p=1000$, $\omega_{\rm base}=29999/30000$,
$k_p=2\times10^{-6}$, $k_i=5\times10^{-11}$ and $K_e=0.1$,
\begin{align}
 \gamma&=\frac{25}{1024},&h&=\frac{30000000}{29999},\nonumber\\
 \kappa_I&=\frac3{59998000},&
 K_u&=-\frac{768}{18749375},\nonumber\\
 K_{\Delta e}&=\frac{29999}{7500}
 \simeq3.9998666667.
 \label{detail:eq:fpd_example_gains}
\end{align}
For a one-frame first error, the matched PI and $F_0$ corrections are both
$2.0500016667\times10^{-6}$. Warming up only the Fuzzy previous error
instead gives approximately $5.0001666722\times10^{-8}$. The example
illustrates an initialization mismatch; it is not a new plant simulation.

\subsubsection{Execution-layer conditions}
Identical raw outputs give identical applied outputs under a common
limiter when raw states are retained and anti-windup is disabled.
Tracking or reset requires compatible state changes. For example, after
processing $E_k$, setting
\[
 A^0_k=c_{\rm held},\qquad
 I_k=\frac{c_{\rm held}-k_pE_k}{\kappa_I},\qquad
 e_k=-\gamma E_k
\]
restores the state relation for the next ordinary update. A branch designed
to produce a held output on an immediate extra update requires its own
integrate-first adjustment; the above is not an instruction to add an
extra call. Limiting the internal reference $F_0$ to $[-5,5]$ also changes
the controller, even if the external correction limiter is shared.

Under a deterministic causal plant, the same initial history, measured
inputs, actuation order, external trace and compatible resets imply
closed-loop equality by induction over events. No globally synchronous
physical grid is required for this algebraic equivalence. By contrast,
matching only labels such as ``same gains'' or ``same random seed'' does
not establish these conditions.

\subsection{Deriving the Exact Central Nonlinear Surface}
\label{detail:subsec:fpi_mechanism}

For $0<r\le5$, define
\begin{align}
 \Omega_r&=\{(X,Y):|X|+|Y|\le r\},\nonumber\\
 m&=\min(|X|,|Y|),\qquad s=|X|+|Y|.
 \label{detail:eq:fpi_domain}
\end{align}
The same-sign interaction is
\begin{equation}
 J(X,Y)=\begin{cases}\operatorname{sgn}(X)m,&XY>0,\\0,&XY\le0.\end{cases}
 \label{detail:eq:fpi_interaction}
\end{equation}
For points in $\Omega_5$, put $a=|X|/5$, $b=|Y|/5$.
Because $a+b\le1$, each input has only its central and sign-adjacent
membership active. Their four joint min weights are
\begin{align}
 w_{00}&=\min(1-a,1-b)=1-\max(a,b),\nonumber\\
 w_{s0}&=\min(a,1-b)=a,\nonumber\\
 w_{0s}&=\min(1-a,b)=b,\nonumber\\
 w_{ss}&=\min(a,b).
 \label{detail:eq:fpd_four_weights}
\end{align}
The subscripts $s$ denote the input's active sign label, not the scalar
$s=|X|+|Y|$. Zero-input boundaries follow by continuity.

The denominator is therefore
\begin{equation}
 D=1-\max(a,b)+a+b+\min(a,b)
  =1+2\min(a,b)=1+2m/5.
 \label{detail:eq:fpd_denominator}
\end{equation}
The central-central consequent is zero. The sign-central and
central-sign rules contribute $(5/2)\operatorname{sgn}(X)a=X/2$
and $(5/2)\operatorname{sgn}(Y)b=Y/2$.
For equal signs, the joint consequent is
$5\operatorname{sgn}(X)$ and contributes $\operatorname{sgn}(X)m$;
for opposing signs it is zero. The four quadrants therefore give
\begin{equation}
 F(X,Y)=\begin{cases}
 \dfrac{F_0+m}{1+2m/5},&X>0,\ Y>0,\\[4pt]
 \dfrac{F_0-m}{1+2m/5},&X<0,\ Y<0,\\[4pt]
 \dfrac{F_0}{1+2m/5},&X>0,\ Y<0,\\[4pt]
 \dfrac{F_0}{1+2m/5},&X<0,\ Y>0.
 \end{cases}
 \label{detail:eq:fpi_quadrants}
\end{equation}
On either axis $m=0$, giving $F=F_0$. These cases establish the following
identity, including boundaries by continuity.
\begin{btdproposition}[Exact central surface]
\label{detail:prop:fpi_central}
For every $(X,Y)\in\Omega_5$,

\begin{equation}
 \begin{aligned}
 F(X,Y)&=\frac{(X+Y)/2+J(X,Y)}{1+2m/5},\\
 J(X,Y)&=\begin{cases}\operatorname{sgn}(X)m,&XY>0,\\0,&XY\le0.
 \end{cases}
 \end{aligned}
 \label{detail:eq:fpi_central}
\end{equation}
\end{btdproposition}
This is an exact rule-level identity on $\Omega_5$, not a fitted
approximation. Outside this diamond, the cross weights in
\eqref{detail:eq:fpd_four_weights} need not equal $a$ and $b$, so the formula
must not be extended without a new case analysis.

\subsection{Sign-conditioned response, sector and nonlinear geometry}

\subsubsection{Same-sign enhancement}
When $XY>0$, $F_0$ has magnitude $s/2$ and the same sign as $X$.
Equation~\eqref{detail:eq:fpi_central} gives
\begin{equation}
 \frac{F}{F_0}=\frac{1+2m/s}{1+2m/5}.
 \label{detail:eq:fpi_same_sign}
\end{equation}
Since $s\le5$, its numerator is at least the denominator. Since
$m\le s/2$, the numerator is at most 2. Thus $1\le F/F_0\le2$.
Define the normalized nonlinear residual by
$\mathcal R_F(X,Y)=F(X,Y)-F_0(X,Y)$. Subtracting $F_0$ gives
\begin{equation}
 \mathcal R_F=
 \frac{\operatorname{sgn}(X)m(1-s/5)}{1+2m/5},\qquad XY>0.
 \label{detail:eq:fpd_same_residual}
\end{equation}
For instance, $F(0.1,0.1)=0.2/1.04\simeq0.19230769$, compared with
$F_0=0.1$. Both are surface values; the actual increment multiplies by
negative $K_u$.

\subsubsection{Opposing-sign attenuation and zero-output cases}
When $XY<0$, the interaction vanishes, so
\begin{equation}
 F=\frac{F_0}{1+2m/5},\qquad
 \mathcal R_F=-\frac{2m/5}{1+2m/5}F_0.
 \label{detail:eq:fpi_opposite_sign}
\end{equation}
For $(X,Y)=(0.1,-0.05)$, $F_0=0.025$ and
$F=0.025/1.02\simeq0.024509804$. On either axis $m=0$ and $F=F_0$.
On the anti-diagonal $Y=-X$, both outputs are exactly zero; a ratio should
not be evaluated there.

\begin{btdcorollary}[Regional pointwise sector]
\label{detail:cor:fpi_sector}
For $(X,Y)\in\Omega_r$, $0<r\le5$, one can write
\begin{equation}
 F=\alpha(X,Y)F_0,\qquad
 \underline\alpha_r=\frac1{1+r/5}\le\alpha(X,Y)\le2.
 \label{detail:eq:fpi_sector}
\end{equation}
\end{btdcorollary}
\begin{proof}
Same-sign ratios lie in $[1,2]$. For opposing signs with $F_0\ne0$,
$m\le r/2$ gives $1/(1+r/5)\le F/F_0\le1$.
The axes agree with $F_0$. If $F_0=0$, both outputs vanish and any
admissible $\alpha$ preserves the equality.
\end{proof}
The equivalent division-free inequality is
\begin{equation}
 [F-\underline\alpha_r F_0][2F_0-F]\ge0.
 \label{detail:eq:fpi_sector_product}
\end{equation}
In particular, $\Omega_1$ gives $\alpha\in[5/6,2]$, while
$\Omega_5$ gives $[1/2,2]$. These bounds describe pointwise values of
the specific nonlinear map. They are not a gain-scheduled stability proof;
the closed loop must retain the input region and accommodate different
node-dependent, time-varying $\alpha$ values.

\subsubsection{Why an axis match is not a Jacobian linearization}
The same-sign term $J$ is positively homogeneous. For $t>0$ and a fixed
direction $v=(v_1,v_2)$, the central formula yields
\begin{equation}
 F'(0;v)=\lim_{t\downarrow0}\frac{F(tv_1,tv_2)}t
 =\frac{v_1+v_2}{2}+J(v_1,v_2).
 \label{detail:eq:fpd_direction}
\end{equation}
Consequently $F'(0;(1,0))=F'(0;(0,1))=1/2$, but
$F'(0;(1,1))=2$. A linear derivative with the two axis values would
instead give 1 on $(1,1)$, a contradiction. The map is nonlinear and not
Fr\'echet differentiable at the origin, despite exact agreement with
$F_0$ on the coordinate axes. ``Axis-calibrated linear counterpart'' is
therefore the appropriate term.

\subsubsection{Where the central interaction comes from}
As an analytic counterfactual, replace min rule weights by
$w^{\rm prod}_{ab}=\varphi_a(X)\varphi_b(Y)$, retaining the membership functions
and singletons. On $|X|,|Y|\le5$, the active consequents add their two
single-input values and the denominator is
$(\sum_a\varphi_a(X))(\sum_b\varphi_b(Y))=1$. Thus
\begin{equation}
 F_{\rm prod}(X,Y)=X/2+Y/2=F_0(X,Y).
 \label{detail:eq:fpd_product}
\end{equation}
The central nonlinear interaction is tied to normalized min activation,
not merely to linguistic labels. The identity isolates the role of min activation; it does not assert
linearity beyond the central square or an experimental comparison of
product and min inference.

\subsection{Nonlinear memory and what it does not identify}

For identical scaled input pairs $(X_k,Y_k)$, including consistent
previous-error initialization, let $A_k^{\rm F}$ use $F$ and
$A_k^0$ use $F_0$, with equal initial accumulators. Subtracting their
updates gives
\[
 A_k^{\rm F}-A_k^0=(A_{k-1}^{\rm F}-A_{k-1}^0)
                  +K_u\mathcal R_F(X_k,Y_k).
\]
Telescoping yields
\begin{equation}
 A_k^{\rm F}-A_k^0=K_u\sum_{\ell=0}^{k}\mathcal R_F(X_\ell,Y_\ell).
 \label{detail:eq:fpi_memory}
\end{equation}
With zero PI/Fuzzy constructor states and the calibrated parameters,
the same-input bookkeeping can also be written
\begin{equation}
 A_k^{\rm F}=k_pE_k+\kappa_I\sum_{\ell=0}^{k}E_\ell
             +K_u\sum_{\ell=0}^{k}\mathcal R_F(X_\ell,Y_\ell).
 \label{detail:eq:fpd_accounting}
\end{equation}
The first two terms form a PI-like shadow evaluated on the Fuzzy
trajectory's own inputs. The last term records accumulated nonlinear
correction. If later inputs lie on the axes, the new residual is zero
but the accumulated difference can remain nonzero.

This identity requires no accumulator projection/reset. Pure external
clipping with retained raw states leaves the raw identity intact, but the
applied difference must be evaluated through the limiter. Tracking/reset
requires explicit additional state-change terms. Moreover, independently
evolving PI and Fuzzy closed loops generate different measurements and
possibly different physical event schedules. A same-input shadow identity
therefore cannot be relabelled as their observed closed-loop difference.

\subsection{Interpretation and the bridge to the stability section}

The derivation provides implementation-consistent PI calibration,
the exact central nonlinear map, and a regional sector suitable for a
closed-loop certificate. Section~\ref{detail:sec:btd_reference} uses these properties after stating its
plant, sampling, latency and initialization assumptions and proving that
the sector's input region remains invariant. The present algebra does not
change the continuous-clock Callisto plant into a discrete-time plant;
it describes the controller updates embedded in that execution.

Same-sign scaled inputs indicate alignment of the two input channels.
They do not always mean increasing absolute error: a transition from
$e_{k-1}=-10$ to $e_k=1$ has $e_k(e_k-e_{k-1})>0$ even though the
absolute error decreased. Likewise, opposing-sign attenuation concerns a
single increment relative to $F_0$, not an established network jitter
attenuation property. Inter-node frequency differences, buffer excursion, common
frequency and control activity must be evaluated separately.

